\documentclass[12pt]{article}
\usepackage[utf8]{inputenc}
\usepackage[T1]{fontenc}
\usepackage[margin=1in]{geometry}
\usepackage{setspace}
\usepackage{xcolor}
\usepackage[colorlinks=true, urlcolor=blue!60!black, linkcolor=blue!60!black]{hyperref}
\setlength{\parindent}{2em}
\setlength{\parskip}{1.25em}
\renewcommand{\baselinestretch}{1.0}
\begin{document}
\begin{center}
{\LARGE \textit{United States v. Belmont}, 301 U.S. 324 (1937)\par}
\vspace{0.8em}
{\large Case Brief\par}
\end{center}
\vspace{1em}

\section{Memory Jogger}
The Court unanimously reversed and let the United States collect a New York bank deposit assigned to it by the Soviet Union in the Litvinov Assignment; Justice Sutherland's majority held that a sole executive agreement made incident to recognition needs no Senate consent and overrides contrary state law and policy (\href{https://supreme.justia.com/cases/federal/us/301/324/}{Justia}).

\section{Detailed Case Facts}
Prior to 1918, a Russian corporation named Petrograd Metal Works deposited a sum of money with August Belmont, a private banker operating in New York City as August Belmont \& Co. In 1918, following the Bolshevik Revolution, the Soviet Government issued a nationalization decree dissolving the corporation and expropriating all of its property and assets wherever situated, including the New York bank deposit.

For sixteen years, the United States withheld formal diplomatic recognition from the Soviet regime. On November 16, 1933, however, President Franklin D. Roosevelt formally recognized the Soviet Government and established normal diplomatic relations. Coincident with recognition, Soviet Foreign Commissar Maxim Litvinov and President Roosevelt executed an exchange of diplomatic correspondence---the ``Litvinov Assignment''---without submitting the compact to the Senate for advice and consent. Under this sole executive agreement, the Soviet Government released and assigned to the United States all amounts admitted to be due or found to be due to it from American nationals, including the expropriated Petrograd Metal Works deposit, so that the United States could assemble a fund to settle American citizens' outstanding claims against Russia.

\section{Procedural History}
After August Belmont's death, the United States demanded the deposit from his executors and, upon refusal, sued at law in the U.S. District Court for the Southern District of New York. The District Court dismissed the complaint for failure to state a cause of action, and the Second Circuit affirmed (85 F.2d 542). Both lower courts reasoned that the deposit's situs was in New York and that enforcing the Soviet nationalization decree would give effect to an act of confiscation contrary to New York public policy. The Supreme Court granted certiorari.

\section{Judicial Votes}
The Supreme Court \textbf{reversed} the Second Circuit, \textbf{9--0} in the judgment. The case was argued on March 4, 1937, and decided on May 3, 1937 (\href{https://supreme.justia.com/cases/federal/us/301/324/}{Justia}; \href{https://www.law.cornell.edu/supremecourt/text/301/324}{Cornell LII}):

\begin{itemize}
    \item \textbf{Majority Opinion:} Justice Sutherland delivered the opinion of the Court, joined by Chief Justice Hughes and Justices Van Devanter, McReynolds, Butler, and Roberts.
    \item \textbf{Concurring Opinion:} Justice Stone filed an opinion concurring in the judgment, joined by Justices Brandeis and Cardozo.
    \item \textbf{Dissenting Opinions:} None.
\end{itemize}

\section{Holding}
\textbf{Questions Presented:}

\begin{enumerate}
  \item \textbf{Sole Executive Agreements:} Did the President possess constitutional authority, incident to diplomatic recognition of a foreign government, to enter into the Litvinov Assignment without the advice and consent of two-thirds of the Senate under Article II, Section 2?
  \item \textbf{Preemption of State Public Policy:} Could New York's public policy against enforcing foreign confiscatory decrees bar the United States from recovering a bank deposit located in New York that was assigned to it under that sole executive agreement?
  \item \textbf{Extraterritoriality and Act of State:} Did the Fifth Amendment or American public policy against uncompensated confiscation invalidate the Soviet Government's expropriation of a Russian corporation's property?
\end{enumerate}

\textbf{Holding:} \textbf{Reversed (9--0 in judgment).} Writing for a six-Justice majority, Justice Sutherland held that: (1) diplomatic recognition retroactively validated the Soviet Government's acts from its inception and barred U.S. courts from questioning its taking of its own corporation's property; (2) the President possessed independent constitutional authority, as sole organ of external relations, to execute the Litvinov Assignment as part of diplomatic recognition without Senate participation; and (3) a valid sole executive agreement carries federal supremacy in foreign affairs and overrides any conflicting state laws or public policies, including New York's anti-confiscation policy.

\section{Analysis of Opinions}
\subsection{Majority Opinion (Justice Sutherland)}

\textbf{1. Recognition and the Act-of-State Doctrine.} Taking judicial notice of President Roosevelt's November 16, 1933 recognition of the Soviet Union and exchange of ambassadors, Sutherland invoked settled precedent---including \textit{Underhill v. Hernandez} (1897) and \textit{Oetjen v. Central Leather Co.} (1918)---holding that recognition operates retroactively to validate all acts of the recognized foreign government from its commencement. Because the political branches control foreign relations, American courts may not sit in judgment on the validity of a foreign sovereign's expropriation of its own nationals' property.

\textbf{2. Validity of Sole Executive Agreements.} Sutherland rejected the premise that every international compact requires Senate ratification under the Treaty Clause. Drawing on his theory of inherent national sovereignty in external affairs articulated months earlier in \textit{United States v. Curtiss-Wright Export Corp.} (1936), he reasoned that the President's authority to recognize a foreign government necessarily includes the power to settle mutual claims that stand as obstacles to full diplomatic relations. While a formal ``treaty'' requires two-thirds Senate consent, international compacts such as protocols, \textit{modi vivendi}, postal conventions, and claims-settlement assignments fall within the President's independent competency.

\textbf{3. Supremacy Over State Law and Policy.} Even though Article VI's Supremacy Clause expressly mentions ``Treaties'' rather than executive agreements, Sutherland equated the preemptive force of valid executive agreements with that of treaties. Because external sovereignty belongs exclusively to the national government, state laws and policies cannot obstruct the President's conduct of foreign affairs. In international negotiations and compacts, ``state lines disappear,'' and for those purposes ``the State of New York does not exist.''

\textbf{4. Non-Extraterritoriality and Narrow Remedial Scope.} Finally, Sutherland dismissed Fifth Amendment objections. The U.S. Constitution, laws, and policies have no extraterritorial operation except regarding U.S. citizens; Petrograd Metal Works was a Russian corporation whose redress lay solely with its own government. Belmont's executors were mere stakeholders with no personal interest in the funds beyond protection against double liability, and the Court left open whether any third-party domestic creditors might later assert adverse claims.

\subsection{Opinion Concurring in the Judgment (Justice Stone)}

Stone agreed that the United States stated a valid claim against Belmont's executors, but solely on narrower state-law and conflict-of-laws grounds. He emphasized that a state ordinarily may refuse to enforce foreign confiscatory transfers of property located within its borders---including bank deposits payable on demand in the state---when they offend local public policy. Here, however, New York had no demonstrated public policy allowing a mere stakeholder-debtor to resist payment to a recognized foreign sovereign (or its assignee) when no New York creditors had intervened. Crucially, Stone refused to hold that a sole executive agreement automatically overrides state law: neither the Litvinov Assignment nor President Roosevelt's diplomatic notes purported to alter New York law or grant the United States greater rights than the Soviet Union possessed. By deciding the case on narrow New York law grounds, Stone sought to avoid creating a broad doctrine of unilateral presidential preemption and to preserve third-party creditors' rights to intervene under New York law.

\subsection{Subsequent Developments}

\begin{itemize}

\item \textbf{Expansion in \textit{Pink}:} Five years later, in \href{https://supreme.justia.com/cases/federal/us/315/203/}{\textit{United States v. Pink}, 315 U.S. 203 (1942)}, the Court solidified and expanded \textit{Belmont} over Chief Justice Stone's dissent, holding that the Litvinov Assignment overrode New York's distribution of a Russian insurer's assets even against competing foreign creditors.

\item \textbf{Claims Settlement and Congressional Acquiescence:} In \href{https://supreme.justia.com/cases/federal/us/453/654/}{\textit{Dames \& Moore v. Regan}, 453 U.S. 654 (1981)}, the Court upheld President Carter's Algiers Accords suspending private U.S. claims against Iran, citing \textit{Belmont} and \textit{Pink} but anchoring executive power in longstanding congressional acquiescence under Justice Jackson's \textit{Youngstown} framework rather than pure unilateralism.
  \item \textbf{Foreign-Affairs Preemption:} In \href{https://supreme.justia.com/cases/federal/us/539/396/}{\textit{American Insurance Ass'n v. Garamendi}, 539 U.S. 396 (2003)}, a 5--4 Court relied directly on \textit{Belmont} and \textit{Pink} to strike down a California Holocaust-era insurance disclosure statute because it interfered with the President's voluntary-settlement policy embodied in executive agreements with European governments.
  \item \textbf{Separation-of-Powers Guardrails:} In \href{https://supreme.justia.com/cases/federal/us/552/491/}{\textit{Medell\'in v. Texas}, 552 U.S. 491 (2008)}, the Court sharply limited \textit{Belmont}'s reach, holding that a unilateral presidential memorandum could not compel Texas courts to reopen criminal convictions under a non-self-executing International Court of Justice judgment. Chief Justice Roberts confined \textit{Belmont}, \textit{Pink}, and \textit{Dames \& Moore} to a ``narrow set of circumstances'' involving diplomatic claims settlement backed by pervasive congressional acquiescence. Meanwhile, in \href{https://supreme.justia.com/cases/federal/us/576/1/}{\textit{Zivotofsky v. Kerry}, 576 U.S. 1 (2015)}, the Court cited \textit{Belmont} to confirm that formal diplomatic recognition is an exclusive presidential power under the Reception Clause, while disclaiming \textit{Curtiss-Wright}'s broader ``sole organ'' dicta.

\end{itemize}

\section{Examples: Future Applications}
\begin{enumerate}

\item \textbf{Same-Side (State Law Preempted): Claims Settlement Incident to Recognition.} Following a coup and peace settlement in the Republic of Veridia, the President formally recognizes Veridia's new government and executes a sole executive agreement extinguishing all pending U.S. civil claims against Veridia's state-owned airline in exchange for a \$400 million lump-sum compensation fund administered by a federal commission. A state statute in New York purports to preserve New York residents' state-court tort suits against the airline regardless of federal settlement. The New York statute is preempted under \textit{Belmont}, \textit{Pink}, \textit{Dames \& Moore}, and \textit{Garamendi}. Because the claims-settlement agreement directly accompanies diplomatic recognition and reflects a longstanding, congressionally acquiesced executive practice, it requires no Senate ratification and overrides conflicting state law.

\item \textbf{Same-Side (State Policy Cannot Block Collection): An Assigned Foreign Bank Claim.} On recognizing a new government in Country Z, the President accepts, by an exchange of diplomatic notes, that government's assignment to the United States of claims owed to Z's nationalized central bank by American debtors, so that the United States can fund settlements of Americans' claims against Z. A company in State Q owes the central bank on a loan. When the United States sues to collect, State Q's courts refuse, citing the State's policy against giving effect to foreign confiscations. Under \textit{Belmont}, the assignment is a valid sole executive agreement incident to recognition, and ``the external powers of the United States are to be exercised without regard to state laws or policies''; the United States may collect (\href{https://supreme.justia.com/cases/federal/us/301/324/}{Justia}).

\item \textbf{Opposite-Side (Agreement Not Preemptive): Directing State Criminal Courts.} By contrast, suppose the President signs an unratified bilateral executive agreement with a foreign nation promising that U.S. state criminal courts will automatically vacate state felony convictions whenever local police fail to notify that nation's consulate, and issues a Presidential Memorandum directing state judges to ignore state procedural-default rules. Congress has enacted no implementing statute. Unconstitutional and non-preemptive under \textit{Medell\'in v. Texas} and the limiting principles foreshadowed by Justice Stone's \textit{Belmont} concurrence. Unlike \textit{Belmont}'s narrow diplomatic claims settlement incident to recognition, dictating state criminal procedure lacks both exclusive Article II authority and historic congressional acquiescence; the President cannot unilaterally transform an international undertaking into supreme domestic law binding state courts without Congress or a self-executing Senate-ratified treaty.

\item \textbf{Opposite-Side (Statute Prevails): An Executive Agreement Contrary to a Federal Statute.} The President concludes a sole executive agreement with a trading partner promising to admit its steel free of the import duties that a federal statute requires, and directs customs officials to stop collecting the duties. An importer that competes with the foreign mills sues to compel collection. \textit{Belmont} held only that a valid sole executive agreement overrides \textit{state} law and policy; it did not suggest that the President, acting alone, may displace an Act of Congress. Because the agreement conflicts with a valid federal statute enacted under Congress's power over foreign commerce, the statute controls and the duties must be collected.

\item \textbf{Fence-Sitter (Uncertain): Ending State Tort Suits Without a Compensation Fund.} The President, acting without Senate consent or new legislation, signs an agreement with a long-recognized foreign government that terminates all pending state-court suits by Americans against that government's state-owned chemical company over a factory disaster, in exchange for trade concessions but no compensation fund. The executive can invoke \textit{Belmont}, \textit{Pink}, and \textit{Dames \& Moore}: settling nationals' claims is a long-standing presidential practice, and state law must yield. But the agreement is not incident to recognition, provides no substitute remedy, and rests on whatever congressional acquiescence can be shown, which \textit{Medell\'in} said confines these precedents to a ``narrow set of circumstances.'' Whether it preempts the suits is uncertain.

\end{enumerate}

\section{Critique}
\textbf{Bypassing the Treaty Clause and Supremacy Clause Text:} Critics note that Article II, Section 2 requires two-thirds Senate consent for treaties and Article VI grants supreme-law status only to the Constitution, federal statutes, and ``Treaties.'' Allowing a sole executive agreement to preempt state law threatens to let the President make supreme domestic law unilaterally without bicameralism or Senate supermajority checks.

\textbf{Overbroad ``State Lines Disappear'' Rhetoric:} Justice Stone's concurrence exposed Sutherland's opinion as unnecessary judicial maximalism. Even where federal foreign policy is paramount, states routinely govern local property titles, estates, and banking relationships; \textit{Medell\'in} later vindicated Stone's caution against inferring sweeping domestic preemption from unilateral executive acts.

\textbf{Asymmetry With Federal Statutes:} While \textit{Belmont} held that a sole executive agreement overrides state law, settled doctrine holds that a sole executive agreement cannot supersede a valid federal statute---leaving sole executive agreements in an awkward intermediate tier of the Supremacy Clause hierarchy.

\section{Key Quotations}
\begin{quote}
``There are many such compacts, of which a protocol, a modus vivendi, a postal convention, and agreements like that now under consideration are illustrations.''
\end{quote}
\noindent --- Justice Sutherland, \href{https://www.law.cornell.edu/supremecourt/text/301/324}{301 U.S. at 330--331} (quoted in \href{https://www.law.cornell.edu/supremecourt/text/315/203}{\textit{United States v. Pink}, 315 U.S. at 229}).

\begin{quote}
``Plainly, the external powers of the United States are to be exercised without regard to state laws or policies.''
\end{quote}
\noindent --- Justice Sutherland, \href{https://supreme.justia.com/cases/federal/us/301/324/}{301 U.S. at 331}.

\begin{quote}
``... our Constitution, laws and policies have no extraterritorial operation, unless in respect of our own citizens. ... What another country has done in the way of taking over property of its nationals, and especially of its corporations, is not a matter for judicial consideration here. Such nationals must look to their own government for any redress to which they may be entitled.''
\end{quote}
\noindent --- Justice Sutherland, \href{https://supreme.justia.com/cases/federal/us/301/324/}{301 U.S. at 332} (quoted in \href{https://supreme.justia.com/cases/federal/us/315/203}{\textit{Pink}, 315 U.S. at 228}).

\begin{quote}
``There is nothing in either document to suggest that the United States was to acquire or exert any greater rights than its transferor, or that the President, by mere executive action, purported or intended to alter the laws and policy of any state in which the debtor of an assigned claim might reside, or that the United States, as assignee, is to do more than the Soviet government could have done after diplomatic recognition---that is, collect the claims in conformity with those laws.''
\end{quote}
\noindent --- Justice Stone (concurring), \href{https://www.law.cornell.edu/supremecourt/text/301/324}{301 U.S. at 336--337}.


\end{document}
